% Seminar 3 - Page 1, manuscript dated 24-Oct-1994.
% Fragment for % Seminar 3 - Page 1, manuscript dated 24-Oct-1994.
% Fragment for % Seminar 3 - Page 1, manuscript dated 24-Oct-1994.
% Fragment for % Seminar 3 - Page 1, manuscript dated 24-Oct-1994.
% Fragment for \input{seminars/pages/seminar3-pagina1-en}; no preamble.
% EDITORIAL NOTES:
% - The first exercise revisits the derangements problem from Seminar 2.
%   The handwritten symbol resembles delta(n), whereas the final line
%   appears to use D_n. Both notations are represented below at the
%   points where they occur.
% - The permutation notation, inclusion-exclusion argument, and
%   the displayed final summation are retained in their original order.
% - In Problem 2(a), there is a crossed-out argument inside the initial
%   Hom expression; the readable final expression is Hom(X x Y, Z).
% - The inverse of the currying map is called beta in the proof; beta
%   is also used for the separate correspondence in part (b) on the
%   previous seminar page. This is retained as a local name.
% - Some intermediate arrows and calculations are crossed out; only
%   the clear final equalities are typeset.
% - The manuscript states the correspondence in 2(b), but only proves
%   2(a) on this page; no proof of (b) is supplied here.

\section{Seminar 3}

% \phantomsection
% \subsection*{Seminar 3}
% \addcontentsline{toc}{subsection}{Seminar 3}
\begin{flushright}
\textit{24-Oct-1994}
\end{flushright}

\noindent\fbox{1}\quad A permutation $\sigma\in S_n$ is called a
\emph{derangement} if
\[
  \sigma(i)\neq i,\qquad \forall i=1,\ldots,n.
\]
For $n\in\mathbb N^{*}$, let
\[
 \delta(n)=\bigl|\{\sigma\in S_n\mid\sigma\text{ is a derangement}\}\bigr|.
\]
Compute $\delta(n)$.

\medskip
\noindent Let
\[
 A_i=\{\sigma\in S_n\mid\sigma(i)=i\},
 \qquad |A_i|=(n-1)!.
\]
Then
\[
 \delta(n)=n!-\bigl|\{\sigma\in S_n\mid
                   \sigma\text{ is not a derangement}\}\bigr|.
\]
Denote the set in braces by $B$. Thus
\[
 B=A_1\cup\cdots\cup A_n.
\]
By inclusion-exclusion,
\[
\begin{aligned}
 |B|
 &=|A_1\cup\cdots\cup A_n|\\
 &=\sum_{i=1}^{n}|A_i|
    -\sum_{1\le i<j\le n}|A_i\cap A_j|+\cdots
      +(-1)^{n-1}|A_1\cap\cdots\cap A_n|\\[3pt]
 &=\binom n1(n-1)!-\binom n2(n-2)!+\cdots
      +(-1)^{n-1}\binom nn\,0!\\[3pt]
 &=\sum_{k=1}^{n}(-1)^{k-1}\binom nk(n-k)!.
\end{aligned}
\]
Hence, in the final notation of the manuscript,
\[
 \boxed{D_n=n!-|B|=\sum_{k=0}^{n}(-1)^k\binom nk(n-k)!}.
\]

\bigskip
\noindent\fbox{2}\quad Let $X,Y,Z$ be sets.

\medskip
\noindent\textbf{(a)} The function sets are in correspondence:
\[
 \operatorname{Hom}(X\times Y,Z)
 \;\overset{\alpha}{\underset{\beta}{\rightleftarrows}}\;
 \operatorname{Hom}\bigl(X,\operatorname{Hom}(Y,Z)\bigr),
\]
where
\[
 \alpha(f)(x)(y)=f(x,y).
\]

\noindent\textbf{(b)} Likewise,
\[
 \operatorname{Hom}(X,Y\times Z)
 \cong\operatorname{Hom}(X,Y)\times\operatorname{Hom}(X,Z).
\]

\medskip
\noindent\underline{Proof of (a).}
Define the reverse map by
\[
 \beta(g)(a,b)=g(a)(b),\qquad
 g:X\longrightarrow\operatorname{Hom}(Y,Z).
\]
For every such $g$,
\[
 (\alpha\circ\beta)(g)=g.
\]
Indeed, for all $x\in X$ and $y\in Y$,
\[
 \begin{aligned}
 (\alpha\circ\beta)(g)(x)(y)
 &=\alpha\bigl(\beta(g)\bigr)(x)(y)\\
 &=\beta(g)(x,y)\\
 &=g(x)(y).
 \end{aligned}
\]
On the other hand, for $f:X\times Y\longrightarrow Z$,
\[
 (\beta\circ\alpha)(f)=\operatorname{id}_{\operatorname{Hom}(X\times Y,Z)}(f)=f.
\]
Indeed, for every $x\in X$ and $y\in Y$,
\[
 \begin{aligned}
 \beta\bigl(\alpha(f)\bigr)(x,y)
 &=\alpha(f)(x)(y)\\
 &=f(x,y).
 \end{aligned}
\]
; no preamble.
% EDITORIAL NOTES:
% - The first exercise revisits the derangements problem from Seminar 2.
%   The handwritten symbol resembles delta(n), whereas the final line
%   appears to use D_n. Both notations are represented below at the
%   points where they occur.
% - The permutation notation, inclusion-exclusion argument, and
%   the displayed final summation are retained in their original order.
% - In Problem 2(a), there is a crossed-out argument inside the initial
%   Hom expression; the readable final expression is Hom(X x Y, Z).
% - The inverse of the currying map is called beta in the proof; beta
%   is also used for the separate correspondence in part (b) on the
%   previous seminar page. This is retained as a local name.
% - Some intermediate arrows and calculations are crossed out; only
%   the clear final equalities are typeset.
% - The manuscript states the correspondence in 2(b), but only proves
%   2(a) on this page; no proof of (b) is supplied here.

\section{Seminar 3}

% \phantomsection
% \subsection*{Seminar 3}
% \addcontentsline{toc}{subsection}{Seminar 3}
\begin{flushright}
\textit{24-Oct-1994}
\end{flushright}

\noindent\fbox{1}\quad A permutation $\sigma\in S_n$ is called a
\emph{derangement} if
\[
  \sigma(i)\neq i,\qquad \forall i=1,\ldots,n.
\]
For $n\in\mathbb N^{*}$, let
\[
 \delta(n)=\bigl|\{\sigma\in S_n\mid\sigma\text{ is a derangement}\}\bigr|.
\]
Compute $\delta(n)$.

\medskip
\noindent Let
\[
 A_i=\{\sigma\in S_n\mid\sigma(i)=i\},
 \qquad |A_i|=(n-1)!.
\]
Then
\[
 \delta(n)=n!-\bigl|\{\sigma\in S_n\mid
                   \sigma\text{ is not a derangement}\}\bigr|.
\]
Denote the set in braces by $B$. Thus
\[
 B=A_1\cup\cdots\cup A_n.
\]
By inclusion-exclusion,
\[
\begin{aligned}
 |B|
 &=|A_1\cup\cdots\cup A_n|\\
 &=\sum_{i=1}^{n}|A_i|
    -\sum_{1\le i<j\le n}|A_i\cap A_j|+\cdots
      +(-1)^{n-1}|A_1\cap\cdots\cap A_n|\\[3pt]
 &=\binom n1(n-1)!-\binom n2(n-2)!+\cdots
      +(-1)^{n-1}\binom nn\,0!\\[3pt]
 &=\sum_{k=1}^{n}(-1)^{k-1}\binom nk(n-k)!.
\end{aligned}
\]
Hence, in the final notation of the manuscript,
\[
 \boxed{D_n=n!-|B|=\sum_{k=0}^{n}(-1)^k\binom nk(n-k)!}.
\]

\bigskip
\noindent\fbox{2}\quad Let $X,Y,Z$ be sets.

\medskip
\noindent\textbf{(a)} The function sets are in correspondence:
\[
 \operatorname{Hom}(X\times Y,Z)
 \;\overset{\alpha}{\underset{\beta}{\rightleftarrows}}\;
 \operatorname{Hom}\bigl(X,\operatorname{Hom}(Y,Z)\bigr),
\]
where
\[
 \alpha(f)(x)(y)=f(x,y).
\]

\noindent\textbf{(b)} Likewise,
\[
 \operatorname{Hom}(X,Y\times Z)
 \cong\operatorname{Hom}(X,Y)\times\operatorname{Hom}(X,Z).
\]

\medskip
\noindent\underline{Proof of (a).}
Define the reverse map by
\[
 \beta(g)(a,b)=g(a)(b),\qquad
 g:X\longrightarrow\operatorname{Hom}(Y,Z).
\]
For every such $g$,
\[
 (\alpha\circ\beta)(g)=g.
\]
Indeed, for all $x\in X$ and $y\in Y$,
\[
 \begin{aligned}
 (\alpha\circ\beta)(g)(x)(y)
 &=\alpha\bigl(\beta(g)\bigr)(x)(y)\\
 &=\beta(g)(x,y)\\
 &=g(x)(y).
 \end{aligned}
\]
On the other hand, for $f:X\times Y\longrightarrow Z$,
\[
 (\beta\circ\alpha)(f)=\operatorname{id}_{\operatorname{Hom}(X\times Y,Z)}(f)=f.
\]
Indeed, for every $x\in X$ and $y\in Y$,
\[
 \begin{aligned}
 \beta\bigl(\alpha(f)\bigr)(x,y)
 &=\alpha(f)(x)(y)\\
 &=f(x,y).
 \end{aligned}
\]
; no preamble.
% EDITORIAL NOTES:
% - The first exercise revisits the derangements problem from Seminar 2.
%   The handwritten symbol resembles delta(n), whereas the final line
%   appears to use D_n. Both notations are represented below at the
%   points where they occur.
% - The permutation notation, inclusion-exclusion argument, and
%   the displayed final summation are retained in their original order.
% - In Problem 2(a), there is a crossed-out argument inside the initial
%   Hom expression; the readable final expression is Hom(X x Y, Z).
% - The inverse of the currying map is called beta in the proof; beta
%   is also used for the separate correspondence in part (b) on the
%   previous seminar page. This is retained as a local name.
% - Some intermediate arrows and calculations are crossed out; only
%   the clear final equalities are typeset.
% - The manuscript states the correspondence in 2(b), but only proves
%   2(a) on this page; no proof of (b) is supplied here.

\section{Seminar 3}

% \phantomsection
% \subsection*{Seminar 3}
% \addcontentsline{toc}{subsection}{Seminar 3}
\begin{flushright}
\textit{24-Oct-1994}
\end{flushright}

\noindent\fbox{1}\quad A permutation $\sigma\in S_n$ is called a
\emph{derangement} if
\[
  \sigma(i)\neq i,\qquad \forall i=1,\ldots,n.
\]
For $n\in\mathbb N^{*}$, let
\[
 \delta(n)=\bigl|\{\sigma\in S_n\mid\sigma\text{ is a derangement}\}\bigr|.
\]
Compute $\delta(n)$.

\medskip
\noindent Let
\[
 A_i=\{\sigma\in S_n\mid\sigma(i)=i\},
 \qquad |A_i|=(n-1)!.
\]
Then
\[
 \delta(n)=n!-\bigl|\{\sigma\in S_n\mid
                   \sigma\text{ is not a derangement}\}\bigr|.
\]
Denote the set in braces by $B$. Thus
\[
 B=A_1\cup\cdots\cup A_n.
\]
By inclusion-exclusion,
\[
\begin{aligned}
 |B|
 &=|A_1\cup\cdots\cup A_n|\\
 &=\sum_{i=1}^{n}|A_i|
    -\sum_{1\le i<j\le n}|A_i\cap A_j|+\cdots
      +(-1)^{n-1}|A_1\cap\cdots\cap A_n|\\[3pt]
 &=\binom n1(n-1)!-\binom n2(n-2)!+\cdots
      +(-1)^{n-1}\binom nn\,0!\\[3pt]
 &=\sum_{k=1}^{n}(-1)^{k-1}\binom nk(n-k)!.
\end{aligned}
\]
Hence, in the final notation of the manuscript,
\[
 \boxed{D_n=n!-|B|=\sum_{k=0}^{n}(-1)^k\binom nk(n-k)!}.
\]

\bigskip
\noindent\fbox{2}\quad Let $X,Y,Z$ be sets.

\medskip
\noindent\textbf{(a)} The function sets are in correspondence:
\[
 \operatorname{Hom}(X\times Y,Z)
 \;\overset{\alpha}{\underset{\beta}{\rightleftarrows}}\;
 \operatorname{Hom}\bigl(X,\operatorname{Hom}(Y,Z)\bigr),
\]
where
\[
 \alpha(f)(x)(y)=f(x,y).
\]

\noindent\textbf{(b)} Likewise,
\[
 \operatorname{Hom}(X,Y\times Z)
 \cong\operatorname{Hom}(X,Y)\times\operatorname{Hom}(X,Z).
\]

\medskip
\noindent\underline{Proof of (a).}
Define the reverse map by
\[
 \beta(g)(a,b)=g(a)(b),\qquad
 g:X\longrightarrow\operatorname{Hom}(Y,Z).
\]
For every such $g$,
\[
 (\alpha\circ\beta)(g)=g.
\]
Indeed, for all $x\in X$ and $y\in Y$,
\[
 \begin{aligned}
 (\alpha\circ\beta)(g)(x)(y)
 &=\alpha\bigl(\beta(g)\bigr)(x)(y)\\
 &=\beta(g)(x,y)\\
 &=g(x)(y).
 \end{aligned}
\]
On the other hand, for $f:X\times Y\longrightarrow Z$,
\[
 (\beta\circ\alpha)(f)=\operatorname{id}_{\operatorname{Hom}(X\times Y,Z)}(f)=f.
\]
Indeed, for every $x\in X$ and $y\in Y$,
\[
 \begin{aligned}
 \beta\bigl(\alpha(f)\bigr)(x,y)
 &=\alpha(f)(x)(y)\\
 &=f(x,y).
 \end{aligned}
\]
; no preamble.
% EDITORIAL NOTES:
% - The first exercise revisits the derangements problem from Seminar 2.
%   The handwritten symbol resembles delta(n), whereas the final line
%   appears to use D_n. Both notations are represented below at the
%   points where they occur.
% - The permutation notation, inclusion-exclusion argument, and
%   the displayed final summation are retained in their original order.
% - In Problem 2(a), there is a crossed-out argument inside the initial
%   Hom expression; the readable final expression is Hom(X x Y, Z).
% - The inverse of the currying map is called beta in the proof; beta
%   is also used for the separate correspondence in part (b) on the
%   previous seminar page. This is retained as a local name.
% - Some intermediate arrows and calculations are crossed out; only
%   the clear final equalities are typeset.
% - The manuscript states the correspondence in 2(b), but only proves
%   2(a) on this page; no proof of (b) is supplied here.

\section{Seminar 3}

% \phantomsection
% \subsection*{Seminar 3}
% \addcontentsline{toc}{subsection}{Seminar 3}
\begin{flushright}
\textit{24-Oct-1994}
\end{flushright}

\noindent\fbox{1}\quad A permutation $\sigma\in S_n$ is called a
\emph{derangement} if
\[
  \sigma(i)\neq i,\qquad \forall i=1,\ldots,n.
\]
For $n\in\mathbb N^{*}$, let
\[
 \delta(n)=\bigl|\{\sigma\in S_n\mid\sigma\text{ is a derangement}\}\bigr|.
\]
Compute $\delta(n)$.

\medskip
\noindent Let
\[
 A_i=\{\sigma\in S_n\mid\sigma(i)=i\},
 \qquad |A_i|=(n-1)!.
\]
Then
\[
 \delta(n)=n!-\bigl|\{\sigma\in S_n\mid
                   \sigma\text{ is not a derangement}\}\bigr|.
\]
Denote the set in braces by $B$. Thus
\[
 B=A_1\cup\cdots\cup A_n.
\]
By inclusion-exclusion,
\[
\begin{aligned}
 |B|
 &=|A_1\cup\cdots\cup A_n|\\
 &=\sum_{i=1}^{n}|A_i|
    -\sum_{1\le i<j\le n}|A_i\cap A_j|+\cdots
      +(-1)^{n-1}|A_1\cap\cdots\cap A_n|\\[3pt]
 &=\binom n1(n-1)!-\binom n2(n-2)!+\cdots
      +(-1)^{n-1}\binom nn\,0!\\[3pt]
 &=\sum_{k=1}^{n}(-1)^{k-1}\binom nk(n-k)!.
\end{aligned}
\]
Hence, in the final notation of the manuscript,
\[
 \boxed{D_n=n!-|B|=\sum_{k=0}^{n}(-1)^k\binom nk(n-k)!}.
\]

\bigskip
\noindent\fbox{2}\quad Let $X,Y,Z$ be sets.

\medskip
\noindent\textbf{(a)} The function sets are in correspondence:
\[
 \operatorname{Hom}(X\times Y,Z)
 \;\overset{\alpha}{\underset{\beta}{\rightleftarrows}}\;
 \operatorname{Hom}\bigl(X,\operatorname{Hom}(Y,Z)\bigr),
\]
where
\[
 \alpha(f)(x)(y)=f(x,y).
\]

\noindent\textbf{(b)} Likewise,
\[
 \operatorname{Hom}(X,Y\times Z)
 \cong\operatorname{Hom}(X,Y)\times\operatorname{Hom}(X,Z).
\]

\medskip
\noindent\underline{Proof of (a).}
Define the reverse map by
\[
 \beta(g)(a,b)=g(a)(b),\qquad
 g:X\longrightarrow\operatorname{Hom}(Y,Z).
\]
For every such $g$,
\[
 (\alpha\circ\beta)(g)=g.
\]
Indeed, for all $x\in X$ and $y\in Y$,
\[
 \begin{aligned}
 (\alpha\circ\beta)(g)(x)(y)
 &=\alpha\bigl(\beta(g)\bigr)(x)(y)\\
 &=\beta(g)(x,y)\\
 &=g(x)(y).
 \end{aligned}
\]
On the other hand, for $f:X\times Y\longrightarrow Z$,
\[
 (\beta\circ\alpha)(f)=\operatorname{id}_{\operatorname{Hom}(X\times Y,Z)}(f)=f.
\]
Indeed, for every $x\in X$ and $y\in Y$,
\[
 \begin{aligned}
 \beta\bigl(\alpha(f)\bigr)(x,y)
 &=\alpha(f)(x)(y)\\
 &=f(x,y).
 \end{aligned}
\]
